% --- Python code listings ---
% ----------------------------
%
% Upstream's \lstset is tuned for Scheme in ways that are actively wrong for
% Python: identifiers absorb symbol characters (alsodigit), quotes are letters
% (alsoletter='), ";" starts a comment, keywords are case-insensitive, and "~"
% is the escape character. Each of those is reset below, scoped to our
% environments so Scheme blocks quoted from the original keep rendering as the
% original renders them.

\definecolor{PythonLight}{HTML}{686868}
\definecolor{PythonDark}{HTML}{262626}
\definecolor{PythonBlue}{HTML}{4172A3}
\definecolor{PythonGreen}{HTML}{487818}
\definecolor{PythonBrown}{HTML}{A07040}

\lstdefinestyle{pythonStyle}{%
  language=Python,
  basicstyle=\ttfamily\codeSize,
  % undo the Scheme-specific lexing inherited from upstream's \lstset
  sensitive=true,
  alsodigit={},
  alsoletter={_},
  escapechar={},
  morecomment=[l]{\#},
  morestring=[b]",
  morestring=[b]',
  keywordstyle=\color{PythonBlue},
  commentstyle=\itshape\color{PythonGreen},
  stringstyle=\color{PythonBrown},
  identifierstyle=\color{PythonDark},
  numbers=none,
  frame=l,
  breaklines=true,
  showstringspaces=false,
}

% Inline Python, for prose that names a function or a literal.
\newcommand{\pycode}[1]{\lstinline[style=pythonStyle]{#1}}

% A Python block typed directly into the text. Prefer \pythonFile below.
\lstnewenvironment{pythonBlock}[1][]{%
  \upshape%
  \lstset{style=pythonStyle,#1}%
}{}

% THE preferred way to put Python in the book.
%
% Code included this way lives in a real .py file under code/, which means
% pytest runs it and mypy checks it. Code typed inline into .tex is never
% imported, never executed, never type-checked, and rots silently -- an
% unacceptable failure mode for a book whose stated sequel is about
% correctness.
%
%   \pythonFile{code/ch1/expressions.py}
%   \pythonFile[firstline=12,lastline=20]{code/ch1/expressions.py}
%
\newcommand{\pythonFile}[2][]{%
  \lstinputlisting[style=pythonStyle,#1]{#2}%
}

% --- Syntax skeletons ---
% ------------------------
%
% The general form of a construct -- def <name>(<formal parameters>): -- is
% neither Python nor Scheme but a third thing: grammar notation, in which some
% parts are literal and some are slots to be filled in.
%
% The original writes these inside a listings environment and smuggles the
% metavariables through the escape character (~\(\langle\)~...), which works in
% print and is meaningless anywhere else. Here they get an environment of their
% own, built on alltt rather than listings: alltt keeps the line structure of a
% verbatim block while still expanding macros, so \meta below is an ordinary
% command and no escaping is needed.
\usepackage{alltt}

% A metavariable. Works in prose and inside syntaxForm alike.
\newcommand{\meta}[1]{\ensuremath{\langle}\textsl{#1}\ensuremath{\rangle}}

\newenvironment{syntaxForm}{%
  \par\vspace{0.5\baselineskip}%
  \noindent%
  \begin{alltt}\codeSize%
}{%
  \end{alltt}%
  \vspace{0.3\baselineskip}%
}

% --- breaking long identifiers ---
%
% Upstream writes Scheme's hyphenated names as \code{square\-/list}, where \-/
% is a hyphen the line may break after. Python's names are joined with
% underscores instead, and TeX will not break at one, so a name like
% branch_structure in running prose overflows the measure with no way to fix it
% short of rewording the sentence.
%
% Redefining \code to make each underscore a break opportunity fixes every such
% name at once, and adds only opportunities: a name that already fits is
% unaffected.
\newcommand{\breakableUnderscore}{\textunderscore\allowbreak}
\renewcommand{\code}[1]{{\upshape\ttfamily\let\_\breakableUnderscore #1}}

% --- figures generated by the book's own code ------------------------------
%
% Section 2.2.4's pictures are produced by code/ch2/picture_language.py into
% build/figures/. This includes the PDF for print; tools/pandoc/preprocess.py
% rewrites it to the SVG of the same name for the web edition.
\newcommand{\generatedFigure}[2][\linewidth]{%
  \includegraphics[width=#1]{build/figures/#2.pdf}}

% --- quotations inside footnotes -------------------------------------------
%
% Upstream's biography of William Barton Rogers puts three block quotations
% inside a single footnote. Pandoc cannot parse an environment inside a
% \footnote argument at all -- it stops at the closing brace -- so these are
% marked with an environment of their own, which tools/pandoc/preprocess.py
% flattens to a plain paragraph before pandoc ever sees it.
\newenvironment{quoteInFootnote}
  {\par\smallskip\leftskip=1.2em\rightskip=1.2em\noindent\ignorespaces}
  {\par\smallskip}
